% =========================================================
% Computational Geometry --- Convex Hulls
% =========================================================

\subsection*{Convexity and Hulls}

\begin{topicbox}{Convexity and Convex Hulls}
\begin{exposition}
This formulation makes precise the geometric idea of filling in all straight
segments forced by a point set. Convex combinations give an algebraic language
for the same operation, and the convex hull records the smallest set closed
under that operation.
\end{exposition}

\begin{definitionbox}{Definition}
\begin{definition}[Convex Set]
\label{def:computational-geometry-convex-set}
A set \(S \subseteq \mathbb{R}^2\) is \emph{convex} if for every pair of points
\(a,b\in S\) and every \(\lambda\in[0,1]\),
\[
  \lambda a+(1-\lambda)b\in S.
\]
Equivalently, \(S\) contains the line segment between any two of its points.
\end{definition}
\end{definitionbox}

\begin{remark*}[Standard quantified statement]
The statement applies to the objects named in the hypotheses and asserts the
corresponding containment, extremal, or construction property stated there.
\end{remark*}

\NoLocalDependencies

\begin{definitionbox}{Definition}
\begin{definition}[Convex Combination]
\label{def:computational-geometry-convex-combination}
A \emph{convex combination} of points
\(q_1,\dots,q_n\in\mathbb{R}^2\) is a point of the form
\[
  \sum_{i=1}^{n}\lambda_i q_i,
  \qquad
  \lambda_i\ge 0 \text{ for all } i,
  \qquad
  \sum_{i=1}^{n}\lambda_i=1.
\]
The numbers \(\lambda_i\) are called the weights.
\end{definition}
\end{definitionbox}

\begin{remark*}[Standard quantified statement]
The statement applies to the objects named in the hypotheses and asserts the
corresponding containment, extremal, or construction property stated there.
\end{remark*}

\NoLocalDependencies

\begin{definitionbox}{Definition}
\begin{definition}[Convex Hull]
\label{def:computational-geometry-convex-hull}
The \emph{convex hull} of a set \(P\subseteq\mathbb{R}^2\), written
\(\conv(P)\), is the set of all convex combinations of finitely many points of
\(P\):
\[
  \conv(P)
  =
  \left\{
    \sum_{i=1}^{n}\lambda_i q_i
    \;\middle|\;
    n\ge 1,\ q_i\in P,\ \lambda_i\ge 0,\ \sum_{i=1}^{n}\lambda_i=1
  \right\}.
\]
Equivalently, \(\conv(P)\) is the smallest convex set containing \(P\), i.e. the
intersection of all convex sets that contain \(P\).
\end{definition}
\end{definitionbox}

\begin{remark*}[Standard quantified statement]
The statement applies to the objects named in the hypotheses and asserts the
corresponding containment, extremal, or construction property stated there.
\end{remark*}

\NoLocalDependencies

\begin{lemmabox}{Lemma}
\begin{lemma}[Convex hull as the smallest convex set]
\label{lem:computational-geometry-convex-hull-smallest-convex-set}
Let \(P\subseteq\mathbb{R}^2\). Then \(\conv(P)\) is convex, contains \(P\), and
is contained in every convex set that contains \(P\). Hence
\[
  \conv(P)
  =
  \bigcap\{K\subseteq\mathbb{R}^2 : K\text{ is convex and }P\subseteq K\}.
\]
\end{lemma}
\end{lemmabox}

\begin{remark*}[Proof]
See
\hyperref[prf:computational-geometry-convex-hull-smallest-convex-set]{the proof stub}.
\end{remark*}

\begin{remark*}[Standard quantified statement]
The statement applies to the objects named in the hypotheses and asserts the
corresponding containment, extremal, or construction property stated there.
\end{remark*}

\begin{dependencies}
\begin{itemize}
  \item \hyperref[def:computational-geometry-convex-set]{Convex Set}
  \item \hyperref[def:computational-geometry-convex-combination]{Convex Combination}
  \item \hyperref[def:computational-geometry-convex-hull]{Convex Hull}
\end{itemize}
\end{dependencies}
\end{topicbox}

\subsection*{Supporting Functionals}

\begin{topicbox}{Supporting Functionals and Vertices}
\begin{exposition}
This topic isolates the passage from geometry to optimization. A supporting
line is a geometric certificate that a convex set lies on one side of a
boundary, while a linear functional turns that certificate into a maximum.
\end{exposition}

\begin{definitionbox}{Definition}
\begin{definition}[Extreme Point]
\label{def:computational-geometry-extreme-point}
Let \(S\subseteq\mathbb{R}^2\) be convex. A point \(p\in S\) is an
\emph{extreme point} of \(S\) if whenever
\[
  p=\lambda a+(1-\lambda)b,
  \qquad
  a,b\in S,\quad \lambda\in(0,1),
\]
it follows that \(a=b=p\). For a finite point set \(P\), the extreme points of
\(\conv(P)\) are called the vertices of the convex hull.
\end{definition}
\end{definitionbox}

\begin{remark*}[Standard quantified statement]
The statement applies to the objects named in the hypotheses and asserts the
corresponding containment, extremal, or construction property stated there.
\end{remark*}

\NoLocalDependencies

\begin{definitionbox}{Definition}
\begin{definition}[Supporting Line]
\label{def:computational-geometry-supporting-line}
A line \(\ell\subseteq\mathbb{R}^2\) is a \emph{supporting line} of a convex set
\(S\) at a point \(p\in S\) if \(p\in\ell\) and \(S\) lies entirely in one of the
two closed half-planes bounded by \(\ell\). Equivalently, there exists
\(d\in\mathbb{R}^2\setminus\{0\}\) such that
\[
  \ell=\{x\in\mathbb{R}^2:\langle d,x\rangle=\langle d,p\rangle\}
\]
and
\[
  \langle d,q\rangle\le \langle d,p\rangle
  \qquad
  \text{for all } q\in S.
\]
\end{definition}
\end{definitionbox}

\begin{remark*}[Standard quantified statement]
The statement applies to the objects named in the hypotheses and asserts the
corresponding containment, extremal, or construction property stated there.
\end{remark*}

\NoLocalDependencies

\begin{lemmabox}{Lemma}
\begin{lemma}[Linear functionals preserve convex combinations]
\label{lem:computational-geometry-linear-functionals-convex-combinations}
Let \(d\in\mathbb{R}^2\), and let
\[
  z=\sum_{i=1}^{n}\lambda_i q_i
\]
be a convex combination of points \(q_1,\dots,q_n\in\mathbb{R}^2\). Then
\[
  \langle d,z\rangle
  =
  \sum_{i=1}^{n}\lambda_i\langle d,q_i\rangle.
\]
Consequently, if \(\langle d,q_i\rangle\le M\) for every \(i\), then
\[
  \langle d,z\rangle\le M.
\]
\end{lemma}
\end{lemmabox}\begin{remark*}[Proof]
See
\hyperref[prf:computational-geometry-linear-functionals-convex-combinations]{the proof stub}.
\end{remark*}

\begin{remark*}[Standard quantified statement]
The statement applies to the objects named in the hypotheses and asserts the
corresponding containment, extremal, or construction property stated there.
\end{remark*}

\begin{theorembox}{Theorem}

\begin{dependencies}
\begin{itemize}
  \item \hyperref[def:computational-geometry-convex-combination]{Convex Combination}
  \item \hyperref[def:addition-on-r]{Addition on $\mathbb{R}$}
  \item \hyperref[def:multiplication-on-r]{Multiplication on $\mathbb{R}$}
\end{itemize}
\end{dependencies}
\begin{theorem}[Supporting functional characterization of hull vertices]
\label{thm:computational-geometry-supporting-functional-hull-vertices}
Let \(P\subset\mathbb{R}^2\) be finite and in general position. A point \(p\in P\)
is a vertex of \(\conv(P)\) if and only if there exists a direction
\(d\in\mathbb{R}^2\setminus\{0\}\) such that
\[
  \langle d,p\rangle>\langle d,q\rangle
  \qquad
  \text{for all } q\in P\setminus\{p\}.
\]
Equivalently, \(p\) is the unique point of \(P\) attaining
\[
  \max_{q\in P}\langle d,q\rangle.
\]
Thus every hull vertex is a maximum of \(P\) under some linear functional.
\end{theorem}
\end{theorembox}

\begin{remark*}[Proof]
See
\hyperref[prf:computational-geometry-supporting-functional-hull-vertices]{the proof stub}.
\end{remark*}

\begin{remark*}[Standard quantified statement]
The statement applies to the objects named in the hypotheses and asserts the
corresponding containment, extremal, or construction property stated there.
\end{remark*}

\begin{remark*}[Finite maximum form]
The maximum is used here because \(P\) is finite. This is the finite-set version
of the same supporting-functional intuition that appears in supremum language
for more general sets.
\end{remark*}

\begin{dependencies}
\begin{itemize}
  \item \hyperref[def:computational-geometry-convex-combination]{Convex Combination}
  \item \hyperref[def:computational-geometry-convex-hull]{Convex Hull}
  \item \hyperref[def:computational-geometry-extreme-point]{Extreme Point}
  \item \hyperref[def:computational-geometry-supporting-line]{Supporting Line}
  \item \hyperref[lem:computational-geometry-linear-functionals-convex-combinations]{Linear Functionals Preserve Convex Combinations}
\end{itemize}
\end{dependencies}
\end{topicbox}

\subsection*{Upper Hull Structure}

\begin{topicbox}{Monotone Slopes and Greedy Construction}
\begin{exposition}
This topic converts the supporting-line description into an ordered planar
construction. Under the general-position assumptions, each relevant maximum is
unique, and the upper hull can be recognized by the strict monotonicity of
successive slopes.
\end{exposition}

\begin{assumption}[General Position for the Hull Theorems]
\label{ass:computational-geometry-general-position}
Let \(P \subset \mathbb{R}^2\) be a finite set in general position: the points
have distinct \(x\)-coordinates and no three points are collinear. For
\(p \in \mathbb{R}^2\), write \(x(p)\) and \(y(p)\) for its coordinates.
\end{assumption}

\begin{theorembox}{Theorem}
\begin{theorem}[Monotone slope characterization of the upper hull]
\label{thm:computational-geometry-monotone-slope-upper-hull}
Let \(P\subset\mathbb{R}^2\) be finite and in general position. Let
\[
  v_0,v_1,\dots,v_m
\]
be the vertices of the upper hull of \(P\), listed in increasing
\(x\)-coordinate, and define
\[
  s_k
  =
  \frac{y(v_{k+1})-y(v_k)}{x(v_{k+1})-x(v_k)},
  \qquad
  k=0,1,\dots,m-1.
\]
Then
\[
  s_0>s_1>\cdots>s_{m-1}.
\]
Conversely, if points \(w_0,\dots,w_m\) satisfy
\[
  x(w_0)<x(w_1)<\cdots<x(w_m)
\]
and their consecutive slopes are strictly decreasing, then every \(w_k\) is a
vertex of
\[
  \conv(\{w_0,\dots,w_m\}).
\]
The lower hull is characterized similarly, with strictly increasing consecutive
slopes.
\end{theorem}
\end{theorembox}\begin{remark*}[Proof]
See
\hyperref[prf:computational-geometry-monotone-slope-upper-hull]{the proof stub}.
\end{remark*}

\begin{remark*}[Standard quantified statement]
The statement applies to the objects named in the hypotheses and asserts the
corresponding containment, extremal, or construction property stated there.
\end{remark*}

\begin{theorembox}{Theorem}

\begin{dependencies}
\begin{itemize}
  \item \hyperref[def:computational-geometry-convex-hull]{Convex Hull}
  \item \hyperref[def:computational-geometry-supporting-line]{Supporting Line}
  \item \hyperref[lem:computational-geometry-linear-functionals-convex-combinations]{Linear Functionals Preserve Convex Combinations}
  \item \hyperref[thm:computational-geometry-supporting-functional-hull-vertices]{Supporting Functional Characterization of Hull Vertices}
\end{itemize}
\end{dependencies}
\begin{theorem}[Greedy slope construction of the upper hull]
\label{thm:computational-geometry-greedy-slope-upper-hull}
Let \(P\subset\mathbb{R}^2\) be finite and in general position. Define
\[
  v_0=\argmin_{p\in P}x(p).
\]
Given \(v_k\) with
\[
  x(v_k)<\max_{p\in P}x(p),
\]
define
\[
  v_{k+1}
  =
  \argmax_{\substack{q\in P\\x(q)>x(v_k)}}
  \frac{y(q)-y(v_k)}{x(q)-x(v_k)}.
\]
Terminate once
\[
  v_k=\argmax_{p\in P}x(p).
\]
Then the points
\[
  v_0,v_1,\dots,v_m
\]
produced by this construction are exactly the vertices of the upper hull of
\(P\), and the slopes encountered along the way are strictly decreasing.
\end{theorem}
\end{theorembox}

\begin{remark*}[Proof]
See
\hyperref[prf:computational-geometry-greedy-slope-upper-hull]{the proof stub}.
\end{remark*}

\begin{remark*}[Standard quantified statement]
The statement applies to the objects named in the hypotheses and asserts the
corresponding containment, extremal, or construction property stated there.
\end{remark*}

\begin{remark*}[Uniqueness of choices]
The distinct \(x\)-coordinates make the initial and terminal \(x\)-extreme
points unique. The no-three-collinear condition makes the slope maximizer in
the greedy step unique.
\end{remark*}

\begin{remark*}[Proof roadmap]
The proof ladder is:
\[
\text{convex combinations}
\to
\text{linear functionals}
\to
\text{supporting lines}
\to
\text{monotone slopes}
\to
\text{greedy hull construction}.
\]
The first two lemmas are algebraic. The monotone slope theorem is an
analytic-geometry statement about three ordered points. The greedy construction
theorem is the algorithmic consequence.
\end{remark*}

\begin{dependencies}
\begin{itemize}
  \item \hyperref[def:computational-geometry-convex-hull]{Convex Hull}
  \item \hyperref[def:computational-geometry-supporting-line]{Supporting Line}
  \item \hyperref[lem:computational-geometry-linear-functionals-convex-combinations]{Linear Functionals Preserve Convex Combinations}
  \item \hyperref[thm:computational-geometry-supporting-functional-hull-vertices]{Supporting Functional Characterization of Hull Vertices}
  \item \hyperref[thm:computational-geometry-monotone-slope-upper-hull]{Monotone slope characterization of the upper hull}
\end{itemize}
\end{dependencies}
\end{topicbox}
